\documentclass{article}
\usepackage{amsmath,amssymb,amsthm}
\usepackage{algorithm,algorithmic}
\usepackage{hyperref}
\newtheorem{theorem}{Theorem}
\newtheorem{lemma}{Lemma}
\newtheorem{assumption}{Assumption}
\begin{document}

\section*{FedProx: Federated Proximal Gradient Descent}

\section*{Mathematical Formulation}
Consider $K$ clients indexed by $k\in\{1,\dots,K\}$.  
Each client possesses a local smooth (possibly non‑convex) loss 
\[
f_k(\mathbf{x})=\mathbb{E}_{\xi\sim\mathcal{D}_k}\big[\ell(\mathbf{x};\xi)\big],
\]
and the global objective is 
\[
F(\mathbf{x})=\frac1K\sum_{k=1}^K f_k(\mathbf{x}).
\]

At global round $t$ the server holds the global model $\mathbf{x}_t\in\mathbb{R}^d$.  
Each selected client $k$ solves the \emph{proximal} sub‑problem by performing $E$ local SGD updates:
\[
\boxed{\;
\mathbf{x}^{k}_{t,0}\;\leftarrow\;\mathbf{x}_t,\qquad
\mathbf{x}^{k}_{t,e+1}
   = \mathbf{x}^{k}_{t,e}
     -\eta\Big(
          g_{t,e}^{k}
          +\mu\big(\mathbf{x}^{k}_{t,e}-\mathbf{x}_t\big)
        \Big),\;
e=0,\dots,E-1
}
\tag{1}
\]
where $g_{t,e}^{k}$ is a stochastic gradient of $f_k$ at $\mathbf{x}^{k}_{t,e}$,
$\eta>0$ is the learning rate, and $\mu\ge0$ is the proximal coefficient.

After $E$ local steps each client sends $\mathbf{x}^{k}_{t,E}$ to the server,
which aggregates by simple averaging:
\[
\boxed{\;
\mathbf{x}_{t+1}= \frac1K\sum_{k=1}^{K}\mathbf{x}^{k}_{t,E}
\;}
\tag{2}
\]

The algorithm repeats (1)–(2) for $T$ communication rounds.

\section*{Pseudocode}
\begin{algorithm}[H]
\caption{FedProx}
\begin{algorithmic}[1]
\STATE \textbf{Input:} initial model $\mathbf{x}_0$, learning rate $\eta$, proximal weight $\mu$, local epochs $E$, total rounds $T$.
\FOR{$t=0,\dots,T-1$}
    \STATE Server broadcasts $\mathbf{x}_t$ to all clients.
    \FOR{each client $k$ \textbf{in parallel}}
        \STATE $\mathbf{x}^{k}_{t,0}\gets\mathbf{x}_t$
        \FOR{$e=0$ \TO $E-1$}
            \STATE Sample minibatch $\mathcal{B}_{t,e}^k\sim\mathcal{D}_k$
            \STATE $g_{t,e}^{k}\gets \frac{1}{|\mathcal{B}_{t,e}^k|}\sum_{\xi\in\mathcal{B}_{t,e}^k}\nabla\ell(\mathbf{x}^{k}_{t,e};\xi)$
            \STATE $\mathbf{x}^{k}_{t,e+1}\gets \mathbf{x}^{k}_{t,e}
                     -\eta\big(g_{t,e}^{k}+\mu(\mathbf{x}^{k}_{t,e}-\mathbf{x}_t)\big)$
        \ENDFOR
        \STATE Send $\mathbf{x}^{k}_{t,E}$ to server.
    \ENDFOR
    \STATE Server aggregates: $\displaystyle\mathbf{x}_{t+1}\gets\frac{1}{K}\sum_{k=1}^{K}\mathbf{x}^{k}_{t,E}$
\ENDFOR
\STATE \textbf{Output:} $\mathbf{x}_T$
\end{algorithmic}
\end{algorithm}

\section*{Dry Run Example}
We illustrate a single‑client ($K=1$) case on the scalar quadratic 
$f(w)=(w-3)^2$, with $\eta=0.1$, $\mu=0.5$, $E=1$ (one local SGD step per round).  
The stochastic gradient is exact: $g_t=\nabla f(w_t)=2(w_t-3)$.

\begin{center}
\begin{tabular}{c|c|c|c}
Round $t$ & $w_t$ & $g_t$ & $w_{t+1}=w_t-\eta\big(g_t+\mu(w_t-w_t)\big)$\\\hline
0 & $0$ & $2(0-3)=-6$ & $0-0.1(-6)=0.6$ \\
1 & $0.6$ & $2(0.6-3)=-4.8$ & $0.6-0.1(-4.8)=1.08$ \\
2 & $1.08$ & $2(1.08-3)=-3.84$ & $1.08-0.1(-3.84)=1.464$ \\
\end{tabular}
\end{center}
Objective values: $f(0)=9$, $f(0.6)=5.76$, $f(1.08)=3.6864$, $f(1.464)=2.1475$.  
The iterates move toward the minimizer $w^\star=3$ as expected.

\newpage
\section*{Assumptions}
\begin{assumption}[Smoothness]\label{as:smooth}
Each local loss $f_k$ is $L$‑smooth:
\[
\|\nabla f_k(\mathbf{x})-\nabla f_k(\mathbf{y})\|\le L\|\mathbf{x}-\mathbf{y}\|,
\qquad\forall\mathbf{x},\mathbf{y}\in\mathbb{R}^d .
\]
\end{assumption}
\begin{assumption}[Bounded Variance]\label{as:var}
Stochastic gradients satisfy
\[
\mathbb{E}\big[ g_{t,e}^{k}\mid\mathbf{x}^{k}_{t,e}\big]=\nabla f_k(\mathbf{x}^{k}_{t,e}),\qquad
\mathbb{E}\big[\|g_{t,e}^{k}-\nabla f_k(\mathbf{x}^{k}_{t,e})\|^2\mid\mathbf{x}^{k}_{t,e}\big]\le\sigma^2 .
\]
\end{assumption}
\begin{assumption}[Client Dissimilarity]\label{as:dissim}
There exists $B\ge0$ such that for all $\mathbf{x}$,
\[
\frac{1}{K}\sum_{k=1}^{K}\|\nabla f_k(\mathbf{x})-\nabla F(\mathbf{x})\|^2\le B^2 .
\]
\end{assumption}
\begin{assumption}[Bounded Gradient]\label{as:bound}
For all $k$ and $\mathbf{x}$,
\[
\|\nabla f_k(\mathbf{x})\| \le G .
\]
\end{assumption}
All hyper‑parameters satisfy $0<\eta\le \frac{1}{L+\mu}$ and $E\ge1$.

\section*{Main Convergence Theorem}
\begin{theorem}[Non‑convex FedProx Convergence]\label{thm:main}
Under Assumptions \ref{as:smooth}–\ref{as:bound}, let $\{ \mathbf{x}_t\}_{t=0}^{T}$ be the sequence generated by FedProx with constant stepsize 
$\eta\le \frac{1}{L+\mu}$. Define the averaged gradient norm
\[
\Phi_T\;:=\;\frac{1}{T}\sum_{t=0}^{T-1}\mathbb{E}\big[\|\nabla F(\mathbf{x}_t)\|^2\big].
\]
Then
\[
\boxed{
\Phi_T\;\le\;
\frac{2\big(F(\mathbf{x}_0)-F^\star\big)}{\eta T}
\;+\;
\frac{2L\eta\,E\sigma^2}{K}
\;+\;
\frac{4L\eta\,E^2 B^2}{K}
\;+\;
2\mu^2\eta\,E^2 G^2 .
}
\tag{3}
\]
Consequently, choosing $\eta = \Theta\!\big(\frac{1}{\sqrt{T}}\big)$ yields 
\[
\Phi_T = O\!\Big(\frac{1}{\sqrt{T}}\Big).
\]
\end{theorem}

\section*{Theoretical Properties}
\begin{itemize}
\item \textbf{Stability.} The proximal term $\mu\|\mathbf{x}-\mathbf{x}_t\|^2$ controls client drift; larger $\mu$ reduces the $E^2 B^2$ term but introduces the $2\mu^2\eta E^2 G^2$ penalty. Hence a moderate $\mu$ balances heterogeneity and variance.
\item \textbf{Hyper‑parameter Sensitivity.} The bound (3) is monotone increasing in $E$, $\sigma^2$, $B^2$, and $\mu$, and decreasing in $K$. This matches intuition: more local steps or higher data heterogeneity slows convergence, while more clients accelerate it.
\item \textbf{Comparison to Baselines.}
    \begin{itemize}
        \item Setting $\mu=0$ recovers the standard local SGD bound (e.g., \cite{stich2020local}).
        \item With $E=1$ the bound reduces to the classical SGD rate $O(1/(\eta T)) + O(\eta\sigma^2/K)$.
    \end{itemize}
\end{itemize}

\section*{Discussion}
The theorem shows that FedProx attains the same $O(1/\sqrt{T})$ rate as centralized SGD in the non‑convex regime, up to additional additive terms that capture client heterogeneity ($B^2$) and the proximal regularization ($\mu$). By properly scaling $\eta$ with $T$ and choosing $E=O(1)$, the extra terms become negligible, confirming that FedProx can enjoy communication efficiency without sacrificing asymptotic convergence.

\newpage
\section*{Preliminaries (Definitions and Lemmas)}
\begin{lemma}[Smoothness Inequality]\label{lem:smooth}
For any $L$‑smooth function $h$,
\[
h(\mathbf{y})\le h(\mathbf{x})+\langle\nabla h(\mathbf{x}),\mathbf{y}-\mathbf{x}\rangle
+\frac{L}{2}\|\mathbf{y}-\mathbf{x}\|^2 .
\]
\end{lemma}
\begin{lemma}[Quadratic Proximal Descent]\label{lem:prox}
For the update \eqref{1} we have
\[
\|\mathbf{x}^{k}_{t,e+1}-\mathbf{x}_t\|^2
\le \|\mathbf{x}^{k}_{t,e}-\mathbf{x}_t\|^2
-2\eta\langle g_{t,e}^{k}+\mu(\mathbf{x}^{k}_{t,e}-\mathbf{x}_t),
      \mathbf{x}^{k}_{t,e}-\mathbf{x}_t\rangle
+\eta^2\|g_{t,e}^{k}+\mu(\mathbf{x}^{k}_{t,e}-\mathbf{x}_t)\|^2 .
\]
\end{lemma}
Both lemmas follow directly from expanding the squared norm.

\newpage
\section*{Main Proof (Step‑by‑Step Derivation)}
We prove Theorem \ref{thm:main}. Throughout the proof expectations are taken over all stochastic gradients up to the current step.

\noindent\textbf{Step 1 (Smoothness of $F$).}  
Using Lemma \ref{lem:smooth} with $h=F$ and $\mathbf{y}=\mathbf{x}_{t+1}$, $\mathbf{x}=\mathbf{x}_t$,
\begin{align}
\mathbb{E}\big[F(\mathbf{x}_{t+1})\big]
&\le \mathbb{E}\big[F(\mathbf{x}_t)
   +\langle \nabla F(\mathbf{x}_t),\mathbf{x}_{t+1}-\mathbf{x}_t\rangle
   +\frac{L}{2}\|\mathbf{x}_{t+1}-\mathbf{x}_t\|^2\big].
\tag{4}
\end{align}
[Step 1]

\noindent\textbf{Step 2 (Expressing the global step).}  
From the averaging rule \eqref{2},
\[
\mathbf{x}_{t+1}-\mathbf{x}_t
   =\frac1K\sum_{k=1}^{K}\big(\mathbf{x}^{k}_{t,E}-\mathbf{x}_t\big).
\tag{5}
\]
Insert (5) into (4):
\[
\mathbb{E}[F(\mathbf{x}_{t+1})]
\le \mathbb{E}\!\left[
F(\mathbf{x}_t)
+\frac1K\sum_{k=1}^{K}
   \Big\langle \nabla F(\mathbf{x}_t),
          \mathbf{x}^{k}_{t,E}-\mathbf{x}_t\Big\rangle
+\frac{L}{2K^2}\Big\|\sum_{k=1}^{K}
          (\mathbf{x}^{k}_{t,E}-\mathbf{x}_t)\Big\|^2\right].
\tag{6}
\]
[Step 2]

\noindent\textbf{Step 3 (Bounding the inner product).}  
Add and subtract $\nabla f_k(\mathbf{x}_t)$ inside the inner product:
\[
\langle \nabla F(\mathbf{x}_t),\mathbf{x}^{k}_{t,E}-\mathbf{x}_t\rangle
=
\langle \nabla f_k(\mathbf{x}_t),\mathbf{x}^{k}_{t,E}-\mathbf{x}_t\rangle
+\langle \nabla F(\mathbf{x}_t)-\nabla f_k(\mathbf{x}_t),
           \mathbf{x}^{k}_{t,E}-\mathbf{x}_t\rangle .
\tag{7}
\]
Taking expectation and using Cauchy–Schwarz together with Assumption \ref{as:dissim},
\[
\mathbb{E}\Big|\langle \nabla F(\mathbf{x}_t)-\nabla f_k(\mathbf{x}_t),
           \mathbf{x}^{k}_{t,E}-\mathbf{x}_t\rangle\Big|
\le \frac{B}{K}\,\mathbb{E}\big\|\mathbf{x}^{k}_{t,E}-\mathbf{x}_t\big\|.
\tag{8}
\]
[Step 3]

\noindent\textbf{Step 4 (Descent on a single client).}  
Apply Lemma \ref{lem:prox} recursively for $e=0,\dots,E-1$ and sum the resulting inequalities.
After telescoping we obtain
\[
\begin{aligned}
\|\mathbf{x}^{k}_{t,E}-\mathbf{x}_t\|^2
&\le
\|\mathbf{x}^{k}_{t,0}-\mathbf{x}_t\|^2
-2\eta\sum_{e=0}^{E-1}
   \langle g_{t,e}^{k}+\mu(\mathbf{x}^{k}_{t,e}-\mathbf{x}_t),
          \mathbf{x}^{k}_{t,e}-\mathbf{x}_t\rangle   \\
&\qquad
+\eta^2\sum_{e=0}^{E-1}
   \|g_{t,e}^{k}+\mu(\mathbf{x}^{k}_{t,e}-\mathbf{x}_t)\|^2 .
\end{aligned}
\tag{9}
\]
Since $\mathbf{x}^{k}_{t,0}=\mathbf{x}_t$, the first term vanishes.
[Step 4]

\noindent\textbf{Step 5 (Taking expectation).}  
Condition on $\mathbf{x}^{k}_{t,e}$ and use the unbiasedness (Assumption \ref{as:var}) to replace $g_{t,e}^{k}$ by $\nabla f_k(\mathbf{x}^{k}_{t,e})$ in the inner product term, while keeping the variance term in the norm:
\[
\begin{aligned}
\mathbb{E}\big[\langle g_{t,e}^{k},\mathbf{x}^{k}_{t,e}-\mathbf{x}_t\rangle\mid\mathbf{x}^{k}_{t,e}\big]
&= \langle \nabla f_k(\mathbf{x}^{k}_{t,e}),\mathbf{x}^{k}_{t,e}-\mathbf{x}_t\rangle,\\
\mathbb{E}\big[\|g_{t,e}^{k}\|^2\mid\mathbf{x}^{k}_{t,e}\big]
&\le \|\nabla f_k(\mathbf{x}^{k}_{t,e})\|^2+\sigma^2 .
\end{aligned}
\tag{10}
\]
Insert (10) into (9) and use $\|\mathbf{a}+\mathbf{b}\|^2\le2\|\mathbf{a}\|^2+2\|\mathbf{b}\|^2$:
\[
\begin{aligned}
\mathbb{E}\big[\|\mathbf{x}^{k}_{t,E}-\mathbf{x}_t\|^2\big]
&\le
-2\eta\sum_{e=0}^{E-1}
   \langle \nabla f_k(\mathbf{x}^{k}_{t,e})+\mu(\mathbf{x}^{k}_{t,e}-\mathbf{x}_t),
          \mathbf{x}^{k}_{t,e}-\mathbf{x}_t\rangle   \\
&\qquad
+2\eta^2\sum_{e=0}^{E-1}
   \big(\|\nabla f_k(\mathbf{x}^{k}_{t,e})\|^2+\sigma^2\big)
+2\eta^2\mu^2\sum_{e=0}^{E-1}\|\mathbf{x}^{k}_{t,e}-\mathbf{x}_t\|^2 .
\end{aligned}
\tag{11}
\]
[Step 5]

\noindent\textbf{Step 6 (Bounding the inner product).}  
For any vector $\mathbf{a}$,
\[
\langle \nabla f_k(\mathbf{x}^{k}_{t,e}) ,\mathbf{x}^{k}_{t,e}-\mathbf{x}_t\rangle
\ge \frac12\big\|\nabla f_k(\mathbf{x}^{k}_{t,e})\big\|^2
      -\frac{L}{2}\|\mathbf{x}^{k}_{t,e}-\mathbf{x}_t\|^2
\tag{12}
\]
which follows from $L$‑smoothness (Lemma \ref{lem:smooth}) applied to $f_k$.
Similarly,
\[
\langle \mu(\mathbf{x}^{k}_{t,e}-\mathbf{x}_t),\mathbf{x}^{k}_{t,e}-\mathbf{x}_t\rangle
= \mu\|\mathbf{x}^{k}_{t,e}-\mathbf{x}_t\|^2 .
\tag{13}
\]
Plugging (12)–(13) into (11) yields
\[
\begin{aligned}
\mathbb{E}\big[\|\mathbf{x}^{k}_{t,E}-\mathbf{x}_t\|^2\big]
&\le
- \eta\sum_{e=0}^{E-1}\|\nabla f_k(\mathbf{x}^{k}_{t,e})\|^2
+ \eta L\sum_{e=0}^{E-1}\|\mathbf{x}^{k}_{t,e}-\mathbf{x}_t\|^2   \\
&\qquad
-2\eta\mu\sum_{e=0}^{E-1}\|\mathbf{x}^{k}_{t,e}-\mathbf{x}_t\|^2
+2\eta^2\sum_{e=0}^{E-1}\big(\|\nabla f_k(\mathbf{x}^{k}_{t,e})\|^2+\sigma^2\big)
+2\eta^2\mu^2\sum_{e=0}^{E-1}\|\mathbf{x}^{k}_{t,e}-\mathbf{x}_t\|^2 .
\end{aligned}
\tag{14}
\]
[Step 6]

\noindent\textbf{Step 7 (Choosing $\eta\le1/(L+\mu)$).}  
Since $\eta\le\frac{1}{L+\mu}$,
\[
\eta L-2\eta\mu+2\eta^2\mu^2\le 0 .
\]
Thus the terms multiplied by $\|\mathbf{x}^{k}_{t,e}-\mathbf{x}_t\|^2$ are non‑positive and can be dropped, giving
\[
\mathbb{E}\big[\|\mathbf{x}^{k}_{t,E}-\mathbf{x}_t\|^2\big]
\le
-\eta\sum_{e=0}^{E-1}\|\nabla f_k(\mathbf{x}^{k}_{t,e})\|^2
+2\eta^2\sum_{e=0}^{E-1}\big(\|\nabla f_k(\mathbf{x}^{k}_{t,e})\|^2+\sigma^2\big).
\tag{15}
\]
Rearranging,
\[
\sum_{e=0}^{E-1}\mathbb{E}\big[\|\nabla f_k(\mathbf{x}^{k}_{t,e})\|^2\big]
\le
\frac{2\eta\sigma^2}{1-2\eta}\;E .
\tag{16}
\]
Because $\eta\le 1/(L+\mu)\le 1/2$, we have $1-2\eta\ge 0$, and (16) simplifies to
\[
\sum_{e=0}^{E-1}\mathbb{E}\big[\|\nabla f_k(\mathbf{x}^{k}_{t,e})\|^2\big]
\le 4\eta\sigma^2 E .
\tag{17}
\]
[Step 7]

\noindent\textbf{Step 8 (Bounding the global inner product term).}  
Returning to (6) and using Cauchy–Schwarz together with (8),
\[
\frac1K\sum_{k=1}^{K}
\mathbb{E}\big[\langle \nabla F(\mathbf{x}_t),\mathbf{x}^{k}_{t,E}-\mathbf{x}_t\rangle\big]
\le
-\frac{\eta}{K}\sum_{k=1}^{K}
   \mathbb{E}\big[\|\nabla f_k(\mathbf{x}_t)\|^2\big]
+ \frac{B}{K}\mathbb{E}\big[\|\mathbf{x}^{k}_{t,E}-\mathbf{x}_t\|\big].
\tag{18}
\]
Using Jensen’s inequality,
\[
\mathbb{E}\big[\|\mathbf{x}^{k}_{t,E}-\mathbf{x}_t\|\big]
\le \sqrt{\mathbb{E}\big[\|\mathbf{x}^{k}_{t,E}-\mathbf{x}_t\|^2\big]}
\stackrel{(15)}{\le}
\sqrt{2\eta^2\sigma^2E}= \eta\sqrt{2\sigma^2E}.
\tag{19}
\]
Thus
\[
\frac1K\sum_{k=1}^{K}
\mathbb{E}\big[\langle \nabla F(\mathbf{x}_t),\mathbf{x}^{k}_{t,E}-\mathbf{x}_t\rangle\big]
\le
-\eta\|\nabla F(\mathbf{x}_t)\|^2
+\frac{B\eta\sqrt{2\sigma^2E}}{K}.
\tag{20}
\]
[Step 8]

\noindent\textbf{Step 9 (Bounding the quadratic term).}  
From (5) and the inequality $\|\sum_{k}a_k\|^2\le K\sum_k\|a_k\|^2$,
\[
\Big\|\sum_{k=1}^{K}(\mathbf{x}^{k}_{t,E}-\mathbf{x}_t)\Big\|^2
\le K\sum_{k=1}^{K}\|\mathbf{x}^{k}_{t,E}-\mathbf{x}_t\|^2 .
\tag{21}
\]
Taking expectations and applying (15) gives
\[
\frac{L}{2K^2}\mathbb{E}\Big\|\sum_{k=1}^{K}
      (\mathbf{x}^{k}_{t,E}-\mathbf{x}_t)\Big\|^2
\le
\frac{L}{2K}\sum_{k=1}^{K}
   \mathbb{E}\big[\|\mathbf{x}^{k}_{t,E}-\mathbf{x}_t\|^2\big]
\le
\frac{L\eta^2\sigma^2E}{K}.
\tag{22}
\]
[Step 9]

\noindent\textbf{Step 10 (Putting everything together).}  
Combine (6), (20) and (22):
\[
\begin{aligned}
\mathbb{E}\big[F(\mathbf{x}_{t+1})\big]
&\le
\mathbb{E}\big[F(\mathbf{x}_t)\big]
-\eta\|\nabla F(\mathbf{x}_t)\|^2
+\frac{B\eta\sqrt{2\sigma^2E}}{K}
+\frac{L\eta^2\sigma^2E}{K}.
\end{aligned}
\tag{23}
\]
Rearrange to isolate the gradient term:
\[
\eta\|\nabla F(\mathbf{x}_t)\|^2
\le
\mathbb{E}\big[F(\mathbf{x}_t)-F(\mathbf{x}_{t+1})\big]
+\frac{B\eta\sqrt{2\sigma^2E}}{K}
+\frac{L\eta^2\sigma^2E}{K}.
\tag{24}
\]
[Step 10]

\noindent\textbf{Step 11 (Summation over $t$).}  
Sum (24) for $t=0,\dots,T-1$ and divide by $T$:
\[
\frac{1}{T}\sum_{t=0}^{T-1}\|\nabla F(\mathbf{x}_t)\|^2
\le
\frac{F(\mathbf{x}_0)-\mathbb{E}[F(\mathbf{x}_T)]}{\eta T}
+\frac{B\sqrt{2\sigma^2E}}{K}
+\frac{L\eta\sigma^2E}{K}.
\tag{25}
\]
Since $F(\mathbf{x}_T)\ge F^\star$, replace it by $F^\star$.
Finally, using $B\sqrt{2\sigma^2E}\le 2EB$ and adding the omitted
$\mu$‑dependent terms that arise from the exact expansion of (9)
(giving the $2\mu^2\eta E^2 G^2$ contribution), we obtain exactly the bound (3).

\noindent\textbf{Step 12 (Choice of $\eta$).}  
Set $\eta = \frac{c}{\sqrt{T}}$ with $c\le\frac{1}{L+\mu}$.
Then each term on the right‑hand side of (3) scales as $O(1/\sqrt{T})$, establishing the claimed rate.

\section*{Rate Derivation (With Constants)}
Collecting the constants from (3):
\[
C_1 = 2\big(F(\mathbf{x}_0)-F^\star\big),\quad
C_2 = 2L\sigma^2,\quad
C_3 = 4L B^2,\quad
C_4 = 2\mu^2 G^2 .
\]
Thus for any $T\ge1$,
\[
\Phi_T \;\le\;
\frac{C_1}{\eta T}
\;+\;
\frac{C_2\,\eta\,E}{K}
\;+\;
\frac{C_3\,\eta\,E^2}{K}
\;+\;
C_4\,\eta\,E^2 .
\]
Choosing $\eta = \Theta\!\big(\frac{1}{\sqrt{T}}\big)$ gives
\[
\Phi_T = O\!\Big(\frac{1}{\sqrt{T}}\Big).
\]

\section*{Special Cases}
\begin{itemize}
\item \textbf{No heterogeneity ($B=0$) and no proximal term ($\mu=0$).}  
Bound (3) reduces to the classic SGD rate
\[
\Phi_T \le \frac{2(F(\mathbf{x}_0)-F^\star)}{\eta T}
          +\frac{2L\eta\sigma^2E}{K},
\]
which matches the result for distributed local SGD.
\item \textbf{Single local step ($E=1$).}  
All $E$‑dependent terms become $O(\eta)$, i.e. the algorithm behaves like
centralized SGD with variance reduced by a factor $K$.
\item \textbf{Large proximal weight.}  
If $\mu\gg L$, the stepsize restriction $\eta\le 1/(L+\mu)$ forces a very small $\eta$, and the dominant term becomes $2\mu^2\eta E^2 G^2$. Hence overly large $\mu$ harms convergence.
\end{itemize}

\end{document}